\section{宣言の三つの責任}
\label{theory:profession:declaration}
2007年のプロフェッション宣言は、自らの技術、顧客、社会に対する責任を示す。能力を高め、正当な守秘を保ちつつ顧客へ誠実に業務を行い、社会と環境への影響を考える\parencite{eth54_profession2007}。顧客だけ、所属企業だけへ責任を狭めない。

倫理綱領は専門職の行動を支える規範である。団体の会員上の適用と、法律・契約・組織規程の適用は別に考える。非会員だから別の義務がすべてなくなるわけではない。反対に、特定団体の内部規範をすべての他団体へそのまま強制できるとも仮定しない。

\section{多様性とDEI：能力を発揮できる条件を整える}
\label{theory:ethics:diversity}
多様性には性別や国籍などの属性だけでなく、知識、経験、専門性、価値観、働き方の違いもある。人の違いを受け入れ、参加の障壁を調べる。違いの尊重は、他者へ危害を与える行動まで無条件に認めることではない。

日本技術士会の2025年宣言は、多様性、公平性、包摂性の三つを結びつける\parencite{eth54_dei2025}。公平性では、一人ひとりの条件に応じて環境や道具を用意する。同じ道具を配るだけでは利用の障壁が残る場合がある。包摂性は、違いを理由に排除せず、能力を発揮して組織の価値を生む状態を目指す。2021年の多様性問題に用いる場合は、後年の補足資料であることを明示する。
